% !PRINTED 204
\weberpage{204}
%<<CONT

has as a consequence, \gesperrt{admits the substitution $S_{\nu,\nu'}$.} But the group
of the $S_{\nu,\nu'}$ is, on account of

\[
S_{\nu,\nu'} S_{\mu,\mu'} = S_{\nu+\mu,\nu'+\mu'}
\]

an Abelian one, which by the formula

\[
S_{\nu,\nu'} = S_{1,0}^{\nu} S_{0,1}^{\nu'}
\]

is representable by the basis $S_{1,0}$, $S_{0,1}$; and so the \gesperrt{general division equation is an Abelian equation and hence algebraically solvable} (Bd. I, § 169\,f.).

\begin{center}
\S\ 61. \quad \gesperrt{The Division of the Periods.}
\end{center}

The problem that still remains to be solved now consists in determining by
algebraic means the quantities

\weq{1}{x_{\mu,\mu'} = \operatorname{sn}\left(\frac{4\mu K + 4\mu' iK'}{n}\right)}

One obtains all values of this quantity by letting $\mu$, $\mu'$, independently
of one another, each run through a complete system of residues
with respect to the modulus $n$. But these values are also
all \gesperrt{different from one another}, for $\operatorname{sn} v$ can be
$= \operatorname{sn} v'$ only if $v'$ is congruent to $v$ or congruent to $2K - v$
modulo $4K$, $2iK'$, and neither of these can occur for two different ones
among the arguments of (1).

The $n^2$ quantities (1) are the roots of the equation

\weq{2}{xA(x^2) = 0,}

and by the formulae

\weq{3}{y = \frac{D(x^2)}{B(x^2)}, \qquad z = \frac{D(x^2)}{C(x^2)}}

one can also express

\[
\begin{aligned}
y_{\mu,\mu'} &= \operatorname{cn}\left(\frac{4\mu K + 4\mu' iK'}{n}\right),\\
z_{\mu,\mu'} &= \operatorname{dn}\left(\frac{4\mu K + 4\mu' iK'}{n}\right)
\end{aligned}
\]

rationally in terms of $x_{\mu,\mu'}$, if the \gesperrt{domain of rationality consists of rational numbers and rational
functions of $x^2$}.

From the roots $x^2$ of the equation $A = 0$ the roots of $B = 0$,
$C = 0$, $D = 0$ can be derived rationally by
§ 60, (2) and § 44, (18), (19), (20); for if $x^2$ is a root
of $A = 0$, then
%CONT>>
